%% psl_paper.tex
%% Semantic-continuity research-agenda edition, October 2026.
%% Earlier manuscripts are preserved under paper/historical/.
%%
%% Compile with: latexmk -pdf psl_paper.tex

\documentclass[11pt,a4paper]{article}

\usepackage{booktabs}
\usepackage{tabularx}
\usepackage{amsmath}
\usepackage{amssymb}
\usepackage[colorlinks=true,linkcolor=blue,citecolor=blue,urlcolor=blue]{hyperref}
\usepackage[T1]{fontenc}
\usepackage[utf8]{inputenc}
\usepackage{geometry}
\geometry{margin=1in}

\begin{document}

\title{Semantic Continuity Across Hardware Generations:\\
       A Falsifiable Research Agenda and Verified Baseline}

\author{Praveen Kumar\\Independent Researcher, Delhi, India}
\date{October 2026}

\maketitle

\begin{abstract}
Long-lived software often outlives the hardware, protocols, and device
representations for which it was first implemented. Existing work already
provides verified compilation, proof-carrying code, platform-independent
models, extensible intermediate representations, legacy emulation, and
verified device drivers. This paper therefore does not claim novelty for any
of those ingredients individually.

PSL investigates a narrower hypothesis: whether a stable semantic contract can
remain unchanged across materially different hardware generations while each
target-specific realization is checked independently, semantically incorrect
or incapable targets fail closed, and unrelated upstream proof evidence
remains reusable.

The current repository provides only a baseline for this research question. A
repaired Lean 4 model proves an axiom-free forward simulation from successful
validation of a legacy 32-bit control representation to successful execution
in its abstract machine. A finite Python-to-Lean bridge currently checks 65
accepted word streams containing 22 distinct ABI words. These results do not
establish hardware-generation continuity. They instead provide a reproducible
starting point from which the stronger realization boundary can be falsified.

The proposed evaluation uses deterministic virtual device models rather than
specialized physical equipment. It requires correct realizations on two
materially different virtual hardware generations, explicit refusal of an
incapable third target, adversarial rejection of wrong bindings,
representations, units, state transitions, and failure handling, and measured
reuse of upstream evidence under hardware substitution.
\end{abstract}

\noindent\textbf{Keywords:}
formal methods, semantic preservation, refinement, hardware evolution,
proof reuse, Lean 4, virtual hardware, verified systems

\section{Motivation}

Software compatibility is often discussed at the level of instruction sets,
binary interfaces, virtual machines, protocol gateways, or source recompilation.
Those mechanisms are valuable, but they do not by themselves answer a semantic
question: when a system is re-targeted, what evidence establishes that the new
realization still performs the intended operation rather than merely producing
well-formed low-level traffic?

A simple example is a setpoint operation. A legacy controller may encode
$25.0^\circ$C as integer 250 in register \texttt{0x20} and require entry into
a configuration mode. A modern controller may use a named parameter, a
floating-point representation, and a different persistence mechanism. Both may
faithfully implement the same intent. Conversely, a driver that writes integer
25 to the legacy register may be syntactically valid while realizing
$2.5^\circ$C.

The working objective is therefore:

\begin{quote}
Software intent should outlive hardware, but only when a target-specific
realization can justify the same observable semantic contract.
\end{quote}

\section{Prior Art and Non-Claims}

Verified compilation already demonstrates strong machine-checked semantic
preservation between source and target programs; CompCert is the canonical
example \cite{leroy2009compcert}. Proof-carrying code established the pattern
of an untrusted producer supplying evidence that a consumer checks
\cite{necula1997pcc}. MLIR provides extensible infrastructure across multiple
levels and heterogeneous targets \cite{lattner2020mlir}. Model Driven
Architecture explicitly separates platform-independent and
platform-specific models \cite{omgMDA}. Formal reasoning about preservation
of legacy computer behavior under emulation also predates this project
\cite{fidge2003legacy}. Recent driver-verification work such as Pancake shows
that realistic low-level device interactions can be subjected to formal
reasoning \cite{pancake2025}.

PSL therefore does not claim to have invented portability, refinement,
certification, hardware abstraction, legacy emulation, or verified drivers.

The research question is whether a useful \emph{continuity boundary} can be
defined so that one semantic contract is independently realized across
hardware generations, erroneous realizations fail closed, and the cost of
hardware evolution is formally localized.

\section{Research Hypotheses}

\subsection{Semantic continuity}

Let $C$ be a semantic contract and $D_1,D_2$ materially different device
models. The first hypothesis is that independent realizations $R_1,R_2$ can
both satisfy $C$ without modifying $C$.

\subsection{Realization soundness}

A declaration that a target supports an operation is not sufficient evidence
that its binding and behavior are correct. A checked relation should reject
realizations with semantically relevant errors such as a wrong resource,
scale, unit, byte order, privilege transition, operation ordering,
persistence property, or failure interpretation.

\subsection{Fail-closed incompatibility}

If a target cannot satisfy the contract, the correct result is refusal. Silent
degradation or hidden emulation is not compatibility unless the contract
explicitly permits the resulting observable behavior.

\subsection{Evolution-local proof reuse}

Replacing $D_1,R_1$ with $D_2,R_2$ should not force re-verification of upstream
obligations that depend only on an unchanged contract and observation model.
This is not assumed; it must be measured and compared with a conventional
refinement baseline.

\section{Candidate Formal Shape}

Let:
\begin{itemize}
\item $C$ denote a semantic contract;
\item $D$ denote a device model;
\item $R$ denote a target-specific realization;
\item $\pi$ denote realization evidence;
\item $Obs$ denote the observation function or relation.
\end{itemize}

The target checker should establish a statement of the general form

\[
  Check(C,D,R,\pi)=\mathsf{true}
  \quad\Longrightarrow\quad
  Obs(Execute(D,R(C))) \models C.
\]

The exact relation may be trace inclusion, forward simulation, refinement, or
another observational relation once nondeterminism, timing, and failures are
made precise.

The important design constraint is that the realization generator itself need
not belong to the trusted computing base.

\section{Current Verified Baseline}

The present repository retains a repaired legacy PVM32 control model because it
provides a concrete checked starting point. The current Lean development
includes:
\begin{itemize}
\item an executable control validator;
\item an abstract machine using the same control transition;
\item an axiom-free theorem that successful word validation implies successful
      abstract execution with the same final control state;
\item a finite Python-to-Lean conformance corpus containing 65 accepted word
      streams and 22 distinct ABI words;
\item fixed adversarial encoded streams that the Lean validator must reject.
\end{itemize}

Lean 4 is pinned and built in CI \cite{moura2021lean4}. Historical QEMU
experiments remain execution evidence rather than formal hardware refinement
\cite{qemu}.

This baseline does not prove the new semantic-continuity hypotheses.

\section{Virtual Evaluation}

The primary experiments should require no specialized equipment.

The first benchmark freezes one semantic contract and defines three
deterministic virtual target generations:

\begin{center}
\begin{tabular}{lll}
\toprule
Target & Characteristic & Expected result \\
\midrule
Legacy A & narrow/fixed-point, stateful protocol & ACCEPT if faithful \\
Modern B & different representation/interface & ACCEPT independently \\
Incapable C & lacks required semantic behavior & REFUSE \\
\bottomrule
\end{tabular}
\end{center}

QEMU, Renode, or a small deterministic emulator may instantiate the executable
device side. The formal claim remains about the stated device model unless a
separate conformance argument connects the executable simulator to that model.

\section{Adversarial Evaluation}

At minimum the evaluation should mutate a correct realization into:

\begin{itemize}
\item wrong target binding;
\item wrong fixed-point scale;
\item wrong unit;
\item wrong byte order or signedness;
\item missing required configuration/authority transition;
\item invalid command ordering;
\item ignored negative acknowledgement;
\item volatile behavior when persistence is required;
\item non-atomic behavior when atomicity is required;
\item a false capability declaration.
\end{itemize}

The significant case is one in which the low-level protocol is syntactically
valid but the semantic realization is wrong.

\section{Proof-Reuse Measurement}

For hardware substitution $D_1\rightarrow D_2$, evaluation should publish:

\begin{itemize}
\item semantic-contract lines changed;
\item upstream semantic proof lines changed;
\item prior certificates invalidated;
\item new device-model obligations;
\item new realization obligations;
\item checker changes;
\item proof-check time and certificate size;
\item percentage of prior evidence reused.
\end{itemize}

A proof-locality claim is meaningful only if those measurements show that
unrelated upstream evidence is actually preserved.

\section{Falsification Criteria}

The research direction should be narrowed or abandoned if:

\begin{enumerate}
\item wrong bindings cannot be rejected without trusting the realization;
\item a second target forces target details into the semantic contract;
\item hardware substitution causes global proof churn;
\item the checker becomes as complex or trusted as the realization;
\item simulator-specific assumptions leak into the semantic theory;
\item a conventional transition-system/refinement formulation achieves the
      same result more directly with no measurable disadvantage.
\end{enumerate}

Negative results are part of the intended research outcome.

\section{Conclusion}

PSL is now a research programme about semantic continuity, not a claim that a
particular syntax, historical vocabulary, or compiler architecture is novel.
The immediate task is to formalize a realization relation strong enough to
distinguish semantically correct behavior from merely well-formed target
operations, and then test whether that relation allows useful proof reuse
across materially different virtual hardware generations.

The ambition remains simple to state:

\begin{quote}
\textbf{Software intent should outlive hardware.}
\end{quote}

\bibliographystyle{plain}
\bibliography{psl_paper}

\end{document}
